\clearpage\section{Convergence of the time integral}
\subsection{Separate the integration cutoff from the vacuum regulator}
At fixed $\epsilon>0$ the early-time integrand is damped. Increasing the upper limit $L$ tests how much of that damped tail was omitted. Reducing $\epsilon$ changes the integral itself, so one must extend the integration far enough to resolve the slower damping. These two tests answer different questions.

At $r=0.03$, the fixed-regulator checks compare $L=80$ to $160$ at $\epsilon=0.1$, $160$ to $320$ at $0.05$, and $320$ to $640$ at $0.025$. The largest absolute change is $2.408\times10^{-16}$. The production calculation keeps $L=640$ and evaluates four regulators:
\begin{center}\begin{tabular}{rrrrr}\toprule $\epsilon$ & 0.1000 & 0.0500 & 0.0250 & 0.0125\\
$10^{10} F_{\epsilon,640}$ & $-2.67977$ & $-3.30831$ & $-3.64204$ & $-3.81267$\\\bottomrule\end{tabular}\end{center}
Fit a quadratic $F_{\epsilon}=a+b\epsilon+c\epsilon^2$ through the three smallest regulators. Evaluating it at zero gives the central value $a$. For the nested nodes $4h,2h,h$ with $h=0.0125$, Lagrange interpolation gives weights at zero
\[
 w_{4h}=\frac{(-2h)(-h)}{(4h-2h)(4h-h)}=\tfrac13,\quad
 w_{2h}=\frac{(-4h)(-h)}{(2h-4h)(2h-h)}=-2,\quad
 w_h=\frac{(-4h)(-2h)}{(h-4h)(h-2h)}=\tfrac83 .
\]
Hence $a=(F_{4h}-6F_{2h}+8F_h)/3=-3.9858101\times10^{-10}$. A least-squares straight line fitted to the same three points gives intercept $-3.9795417\times10^{-10}$. Their absolute difference, $6.2684\times10^{-13}$, measures sensitivity to extrapolation order.

\subsection{Initial state and quadrature checks}
The two fields already mix during the quadratic evolution. Initialising them as an uncorrelated free-field product at finite $X$ can introduce a spurious phase into the residual. The production state matches the coupled positive-frequency solution through sixth order in $1/X$. The starting-time comparison at otherwise fixed settings is
\begin{center}\begin{tabular}{lrr}\toprule Initial condition & $10^{10}F(X=800)$ & $10^{10}F(X=1600)$\\\midrule
Free product & $-3.04257$ & $-4.81085$\\
Sixth-order adiabatic & $-3.9858100903$ & $-3.9858100896$\\\bottomrule\end{tabular}\end{center}
The central run uses $X=1600$. The mode solver runs down to $x=10^{-6}$ so it can evaluate long legs at $rz_f$. The quadrature uses a logarithmic coordinate below $z=1$ (including $\dd z=z\,\dd\ln z$) and linear, oscillation-resolving panels above $1$. Doubling the panel density and raising the Gauss order from 12 to 16 tests discretisation. An independent adaptive Gauss--Kronrod integral at $r=0.03$ and $\epsilon=0.025$ differs from the baseline fixed-node value by $1.68\times10^{-22}$.

\clearpage\subsection{How much of the error matters for the clock?}
Vary one setting at a time and take the absolute change in the extrapolated result. The regulator entry compares the two intercepts. The tail entry uses the largest fixed-regulator cutoff change. The final-time entry is the larger displacement when $z_f=0.01$ is changed to $0.003$ or $0.001$. At $r=0.03$ the six contributions are
\begin{center}\begin{tabular}{lr@{\qquad}lr}\toprule Source & $|\delta F|$ & Source & $|\delta F|$\\\midrule
Initial $X$ & $6.38\times10^{-20}$ & Quadrature & $9.47\times10^{-22}$\\
Regulator & $6.27\times10^{-13}$ & Final time & $4.77\times10^{-13}$\\
ODE tolerances & $3.58\times10^{-19}$ & Early tail & $2.41\times10^{-16}$\\\bottomrule\end{tabular}\end{center}
The ODE check loosens the relative and absolute tolerances from $(10^{-12},10^{-15})$ to $(10^{-11},10^{-14})$. Summing the six absolute changes gives $\delta F_{\rm num}=1.1043\times10^{-12}$, about $0.277\%$ of $|F|$. This sum measures sensitivity to the tested choices, with no assumption of independent random errors.
\fig{bispectrum_convergence}{.78}{Six pilot triangles with the sum of the tested numerical variations. A small total signal near a zero crossing makes pointwise percentages a poor measure of clock resolution.}
At $r=0.001$ the pointwise error rises to $4.77\%$ because $F$ is near zero. Before fitting, the pilot uses a scale estimate from the full signal,
\[
 R_{\rm pilot}=\frac{\max_i[\delta F_i/r_i^{3/2}]}{\max_i[|F_i|/r_i^{3/2}]}=0.005845.
\]
This passes the $0.01$ threshold for a denser scan. The denominator still contains the smooth background. The next section compares the numerical errors with the amplitude of the fitted oscillation itself.
\source{8}{Error budget, cutoff pairs, regulated variations and independent quadrature}
